\section{Hexad}

The Hexad\footnote{\url{http://www.scribd.com/doc/48888397/6/On-the-Hexad}} is the first perfect number: it arises, by multiplication (rather than addition) from the Dyad and the Triad, and is hence termed ``marriage'' (whereas the Pentad is androgynous). That is, the Hexad unifies Male-Female through ``blending'' and harmony, rather than sticking them together through addition in the Pentad, in which case (because the Pentad is 2:3), one or the other prevails; with the Hexad, there is a perfect mixing.

There are six directions: back and forth, up and down, left and right; so Reality's spatial structure reflects the Hexad. Furthermore, if one divides all number series into Triads, they reflect (in one way or another) the Hexad. $1+2+3=6$, $4+5+6 = 15$, and $1+5 = 6$ again (not to mention that the Hexad is a second Triad), $7+8+9= 24$ and $2+4=6$, etc. (Readers who are familiar with traditional astrology will recognize the summing of the individual digits of sums to reduce the number to a simple one).

6 is the sum of it's multiplicative factors, which are all prime ($1\times 2\times 3=6$).

Here are some of its titles\footnote{\url{http://www.american-buddha.com/cult.pythagsourcebook.2.32.htm}}: Resembling Justice The Thunder-Stone, Amphitrite (Poseidon's wife; a verbal pun: on both sides amphis three trias ), Male-Female Marriage, Finest of All In Two Measures, Form of Forms, Peace, Far-Shooting (name of Apollo) Thaleia Kosmos, Possessing Wholeness, Cure-All (panacea), Perfection, Three-Fold Health, Reconciling.

The Pythagorean theory of musical harmonies\footnote{\url{http://www.american-buddha.com/cult.pythagsourcebook.2.33.htm}} (which is really, as Iamblichus states, no ``theory'' at all, but rather ``the way things actually are'') is fairly complex and is also related to Number.

The Pentad is the beginning of Life on the lower planes, \emph{but it is through the Hexad that actual re-integration occurs}. In searching for other insights to this piece, I came across St. Nicholas of Flue\footnote{\url{http://paulijungunusmundus.eu/rfr/radbilde.htm}}, whose had a vision of ``Two Trinities'' : one was of the ``old gods'' of the West, and one was the newer, Christian perichoretic Trinity. Tomberg gives the key to interpreting some of this, when he notes that the two triangles of the key of Solomon symbolize the movement of lower Creation upward and the the movement of Spirit downward, integrated through the mediation of the regenerated man. Thus, we can see that the ``old gods of the West'' are not to be seen as ``un-Christian'', but are rather, to be baptized and lifted up, through a process of cultural acclimation, in which Western Man is re-integrated into the prime Tradition through participation in Logos. This would (obviously) occur at both an individual and a corporate level, and beautifully illustrates how Christendom was designed to ``give the pattern'' or archetype to ``men of the West''.

Nicholas of Flue\footnote{\url{https://en.wikipedia.org/wiki/Nicholas\_of\_Flüe}}, incidentally, is also credited with preserving peace\footnote{\url{http://www.ewtn.com/library/MARY/NICHFLU.HTM}} between the Catholic and Protestant cantons of Switzerland through his sanctity, and of being the father of the ``state'' of Switzerland. The reader can see through this how Iamblichus' ``Numbers'' are not mere arithmetic or private meditation only, but are indeed true understandings of how the entire Cosmos functions, at both lower levels and higher ones, physical or spiritual, and by implication, everything in between as well.


\flrightit{Posted on 2013-02-08 by Logres}

\begin{center}* * *\end{center}

\begin{footnotesize}
\begin{sffamily}
\texttt{David on 2013-02-09 at 10:49 said:}

This series about numbers are really interesting. I will have to read more of Iamblichus. Thank you.

\hfill

\texttt{Logres on 2013-02-10 at 22:33 said:}

Thank you -- I found it interesting to (the original text). What is fascinating is the figure of Nicholas of Flue: he seems to have been a mystic of a different color (from even the traditional picture of one). But then again, maybe they all are. I do think the ``double Trinity''/Hexad configuration is worth pursuing.

\hfill

\texttt{David on 2013-02-10 at 22:43 said:}

And it does raise the question to what degress Pythagoreas can be linked with all that ? He did had alot to say about mysteries and mathematics.

\hfill

\texttt{Logres on 2013-02-11 at 09:06 said:}

I think if you read Iamblichus' Life of Pythagoras, you would get a very good idea of how much: Taylor's translation is supposed to be classic. Iamblichus was contemporary with the Christian revelation, so I am guessing that his exposition (from my perspective) contained a fuller exposition of Tradition than Pythagoras, who had a different task -- that of laying the paternal foundations and groundwork for the creation of Mediterranean Man (as some have called it). That said, Pythagoras is central to the exposition of both Number and Tradition: whatever can be said about Iamblichus' esoteric teachings, can undoubtedly be said to be at least latent in Pythagoras. I think it's impossible to understate how high a degree of veneration the ancients held him in.

\hfill

\texttt{Logres on 2013-02-11 at 13:26 said:}

From the article on Flue:

``Anahata, the suksma aspect of the heart, thus forms the center of two opposed triads. The lower three chakras obviously contain the suksma aspect of what I have called the drive triad, i.e., in psychological language that aspect of exploration, sexuality, and aggression that can be experienced in introversion; the three upper chakras contain the renewed upper Trinity of Logos, meditation, and Eros, so that anahata ultimately links the lower and the upper trinities. It follows that the heart chakra also contains the Seal of Solomon (cf. figure 5), the union of the instinctive with the spiritual/psychic Trinity which is found in the human heart. This union of the double triadic symbol with the heart, however, corresponds to the structure of the wheel image.''

This seems to fit very well with Tomberg's Meditations on the Tarot.

From physics:
\url{http://en.wikipedia.org/wiki/Quark\_model}


\hfill

\texttt{Pieter on 2013-02-13 at 23:03 said:}

I just stumbled on this site and this post, and although I didn't follow all of your points (e.g. the pentad is the beginning of life on the lower planes) I enjoyed the study. I was curious if you or your audience are familiar with Nassim Haramein's work on the geometry of the ``vacuum'' (i.e. nature). Nassim posits that the basic geometrical unit of the vacuum is the star tetrahedron. The hexahedron also plays a very important role in nature as it is the most stable shape in terms of vector equilibrium. Hexads keep popping up over and over in nature.

This is the quickest way to become familiar with Nassim's theory:

\url{https://www.youtube.com/watch?v=S1JDMToJDe0}


The best lecture I'd recommend is the 8-hour Rogue Valley Metaphysical Library (you gents seem to have longer attention spans than that of average Americans, so here you go):

\url{https://www.youtube.com/watch?v=79\_HwQ-92f8}


Going back to the pentad being the beginning of life on the lower plane, I wonder what your meaning is? Perhaps something esoteric? Even the structure of cell division (material plane) occurs by multiplication of two.

And spheres (cells) group in 3s. Imagine pushing 2 marbles of equal radius together. Now add a third marble and you have greater equilibrium with a sort of isosceles triangle. When you add a fourth, then you must set this on top of the the triangular base in order to maintain equilibrium. Otherwise you're forced to make an awkward diamond, which is essentially two isosceles triangles sharing a common center. Spheres group structurally in 3s, not 2s, 4s, 5s, etc.

Where does the pentad fit in?

\hfill

\texttt{Caleb Cooper on 2013-02-14 at 12:26 said:}

Hello Pieter,

A quick (non)answer is your intuition of esotericism is correct. In fact, the practical experience of the esoteric is pretty much the focus of Gornahoor. This isn't to say that looking for links between physics and metaphysics isn't worthwhile, and probably is of interest to some readers here, but it is outside the focus that our host Cologero tries to maintain.

I am sympathetic to the interest in physics though, as it was a gateway for me to Tradition. So hopefully I'll be forgiven by our gracious host for using its profane language as a bridge to try a give a too rough sketch of what's going on here.

My bridge was realizing that Ockham's razor meant supposing a Void (zero) ``in the beginning.'' In quantum mechanics, as space decreases, the production of virtual particles inversely increases, so in a Void where there is no space, there would be an explosion producing an infinite wavestate that would extend above all space and time, containing all possibilities within itself. This could be identified as the monad (1) (the Greeks didn't have zero to work with though, so you can apply a lot of the attributes of the monad to the Void, in which case this is the dyad, 2), or what humans often call God.

A human body is more or less a radio, while the soul is analogous to the music it plays; having it's origin outside the radio and best understood in terms of frequencies and wavestates. The material universe we inhabit is created by wavestate collapses from the original infinite wavestate; that is the spirit of God descends, through multiple non-physical realms that correspond to certain frequencies, and the pentad corresponds to the plane we psychically inhabit.

Gornahoor, and Tradition, are especially interested in the opposite movement; attuning our radio to the higher frequencies so we can ascend back to God. Hope this helps.

\hfill

\texttt{Logres on 2013-02-14 at 17:41 said:}

Caleb, great answer.

Pieter, I wouldn't add much -- I am not a physicist, and have only a lettered man's acquaintance with it, although I am interested in it. Any subject can be a window into Tradition (music, mathematics, etc.): we will be doing some posts on this later on (the Trivium and Quadrivium).

I am sure you are familiar with the idea that certain things are so nebulous that to observe them is to alter them. We can even go further than that and say, to discuss them or think about them, is to alter them, because Tradition teaches that thoughts are not arbitrary ``things'' contained in our heads. Thus, empirical observation and mathematical models based on peering into the Reality fractal can only yield approximated models which are more or less satisfactory for various purposes (although I think Polkinghorne is right -- they have verisimilitude, at their level, and constitute a ``real'' knowledge). This limitation means that all empirical knowledge whatsoever ultimately cannot tell us about the Absolute, which is experienced by intuition and according to various modes which fit the receptor. A brilliant physicist who was also an esoteric student could, no doubt, relate the two and find his actual knowledge of God enhanced by his studies -- I find this to be the case in my own semi-speciality, which is letters (history and literature). All of these studies have a tendency to assume command of the intuitive cognition which makes knowledge actualized, by metaphysically identifying knower, known, and mode of knowing.

As long as this is born in mind, one can approach Physics, and use it to approach Traditional studies. In answer to your question, I will take a stab and say that the Pentagram is the actual commencement of biological arrangement and chemical life. The fact that Hexads are more common in Nature simply means that Nature doesn't flourish into something more than biological soup (ie., it doesn't approach beauty and fullness of being) without ``6''. There may be traces of this left behind in Nature's molecular structure, and in fact are (eg., snowflakes). However, I don't have a concrete empirical answer for why Pentagrams are so absent, because I don't know enough physics to attempt it -- although their absence may be suggestive. Perhaps someone like yourself, if immersed in Tradition, could discover why this is so.


\hfill

\end{sffamily}
\end{footnotesize}

